\documentclass[nogap]{sfs-2487-2024}

% Muistio tiiviillä, sisennetyllä kappaletyylillä. Vaihtoehto [nogap] erottaa
% peräkkäiset kappaleet ensimmäisen rivin sisennyksellä kappalevälin sijaan
% (oletustyylissä kappaleet erotetaan toisistaan kappalevälillä, 6.4.3).

\logo{\includegraphics{logo-organisaatio}}
\doctype{Muistio}
\date{3.6.2024}
\author{Sari Salonen}

\subject{Henkilöstöhallinto}
\title{Etätyökäytäntöjen päivitys}

\contactinfo{Organisaatio Oy}{%
  Katuosoite\\
  Postinumero Postitoimipaikka\\
  hr@organisaatio.fi}

\begin{document}

\maketitle

\marginlabel{Vastaanottajat}

Esihenkilöt ja henkilöstö

\section{Taustaa}

Nykyiset etätyökäytännöt laadittiin poikkeusolojen aikana, ja ne ovat
olleet voimassa muuttumattomina kahden vuoden ajan. Käytäntöjen
toimivuudesta on kertynyt runsaasti kokemusta sekä esihenkilöiltä että
henkilöstöltä.

Saadun palautteen perusteella käytännöt kaipaavat täsmennyksiä erityisesti
läsnäolon ja tavoitettavuuden osalta. Tähän muistioon on koottu keskeiset
muutosehdotukset jatkokäsittelyä varten.

\section{Ehdotetut muutokset}

Etätyötä voi jatkossa tehdä enintään kolmena päivänä viikossa. Vähintään
kaksi päivää tehdään lähtökohtaisesti työpaikalla, jotta tiimien yhteistyö
ja uusien työntekijöiden perehdytys säilyvät sujuvina.

Yhteiset kokouspäivät sovitaan tiimeittäin, ja niistä ilmoitetaan
kalenterissa hyvissä ajoin. Tavoitettavuus työaikana koskee yhtä lailla
etä- ja lähityötä.

\section{Aikataulu}

Päivitetyt käytännöt otetaan käyttöön seuraavan vuosineljänneksen alusta.
Esihenkilöt käyvät muutokset läpi tiimeissään ennen niiden voimaantuloa.

\forinformation{Johtoryhmä}

\makecontactinfo

\end{document}
